% SPDX-License-Identifier: Apache-2.0
% covers: FrameOut
\datasheet{FrameOut}{Words from a fetch engine, bytes to the Ethernet
transmitter, the last of them the frame's last.}

\dsfacts{\code{lib/parts/src/ethdma.rs}}{\code{txhdl\_parts::ethdma::FrameOut}}%
{\code{//docs:examples}}

\subsection*{Function}

\code{LineFetch} moves 32-bit words, because that is what the bus
moves. A frame is a count of bytes, from 64 to 1518 of them, so its
last word holds one, two or three real bytes as often as four.
\code{FrameOut} is that difference on the way out, and nothing else:
it does not touch the bus and does not know where the frame is in
memory.

It takes a word at a time on its input channel and emits its bytes,
lowest first, as \code{EthByte}s. It marks the last byte of the
\emph{frame}, not the last byte of the last word, so the bytes that
sit above the frame's end in a partial last word are never sent.

\code{bytes} is the frame's length and \code{go} asks for a frame.
\code{running} is high while one is going. \code{nwords} is the word
count a fetch engine needs for a frame of \code{bytes}, that is
\code{(bytes + 3) >> 2}.

\subsection*{Parameters}

None. The byte count is 16 bits wide and the words are 32.

\dstables{FrameOut}

\subsection*{Behaviour}

\begin{itemize}
\item \code{go} is a level, not a pulse. A request is taken when the
unit is idle, and \code{bytes} is latched then, so the register block
may hold \code{go} until \code{running} says it was taken.
\item A run ends when the bytes emitted reach the count latched at
the start. A count of zero ends the run as it starts and sends
nothing.
\item A word is taken when the last one is spent, which leaves one
cycle in five with no byte. That costs nothing on the wire: the
transmitter stores a whole frame before it sends any of it, so this
side is never what the line waits for.
\item \code{nwords} is driven from \code{bytes} as it is, not from the
count latched at the start, because the fetch engine reads it in the
cycle it is told to go. The division is here because a board is only
wires between its units, and this is the one unit that stands between
a count of bytes and a count of words.
\end{itemize}

\subsection*{Verification}

\begin{itemize}
\item \code{ex\_ethdma} sends two frames, of 101 and 67 bytes, from
\code{FrameOut} through the real \code{EthTx} and \code{EthRx} to
\code{FrameIn}, and checks every word that comes back. Neither length
divides by four, and the two leave one and three real bytes in their
last words.
\item Two frames, because one cannot catch a byte side that reads past
the count. Such a unit still marks the right byte as the last, so the
first frame is the right length and passes every check. The bytes it
read past that frame wait in the transmitter, which saw no last byte
for them, and go out in front of the next frame.
\item The netlist is checked against that run under nvc and
Verilator: \code{//docs:ethdma\_sim\_frame\_out\_frame\_out\_tb\_test}
and \code{//docs:ethdma\_vsim\_frame\_out\_test}.
\item \code{ex\_framelen} runs it again in front of \code{FrameLen},
with the same two frames.
\item On the board,
\code{two\_frames\_written\_to\_memory\_leave\_the\_port\_as\_written\_in\_order}
in \code{cpu/vreteno/tests/board.rs} has the core write two frames of
23 and 18 bytes and ask for both back to back, and it checks the bytes
that leave the port.
\end{itemize}

\subsection*{Used by}

\code{Board}, as \code{fout}, between the fetch engine \code{fetch}
and \code{EthShare}. \code{EthSlots} gives it \code{bytes} and
\code{go}, and reads \code{running} as \code{tx\_busy}. The examples
\code{ex\_ethdma} and \code{ex\_framelen}.

\subsection*{Limits}

It trusts the count. A count larger than the frame in memory sends
whatever the fetch engine read past it, and nothing here can tell.

It runs on one clock, \code{DefaultClock}.
